\documentclass[11pt]{article}
\usepackage[a4paper,margin=2.6cm]{geometry}
\usepackage{amsmath,amssymb,amsthm}
\usepackage{booktabs}
\usepackage{microtype}
\usepackage[hidelinks]{hyperref}

\newtheorem{theorem}{Theorem}
\newtheorem{proposition}[theorem]{Proposition}
\newtheorem{corollary}[theorem]{Corollary}
\theoremstyle{definition}
\newtheorem{definition}[theorem]{Definition}
\newtheorem{assumption}{Assumption}
\renewcommand{\theassumption}{\Alph{assumption}}
\theoremstyle{remark}
\newtheorem{remark}[theorem]{Remark}

\newcommand{\R}{\mathbb{R}}
\newcommand{\Dom}{\Omega}

\title{An upper bound on the instances added by the real-profile repair}
\author{FormalVerificationRandomForests.jl, branch \texttt{experiments}}
\date{23 October 2026}

\begin{document}
\maketitle

\begin{abstract}
The certification loop retrains a model, asks an SMT solver for a counterexample to the
constraint, and adds that counterexample to the training set.
In the real-profile repair, the features outside the left constraint may only take the
values of a real training instance.
We prove that, if the retrained model keeps every already repaired profile correct
(Assumption~\ref{ass:persistence}), the loop certifies the model after at most $N$
cycles, where $N$ is the number of distinct real profiles, hence after adding at most
$mN$ instances, $m$ being the number of instances added per cycle.
We then show that the assumption cannot be dropped: without it, even with a finite set of
profiles and an exact solver, no finite bound exists.
Finally, we explain why no analogous bound exists for the original repair, and compare the
bound with the breastcancer experiments, where the assumption held exactly in 7 of 15 runs
and the loop always stopped far below the bound.
\end{abstract}

\section{Setting}

Let $d$ be the number of features and let $D_0=\{(x_i,y_i)\}_{i=1}^{n}$, with
$x_i\in\R^d$, be the training set.
The left constraint $\varphi$ mentions a set of feature indexes $K\subseteq\{1,\dots,d\}$
(the \emph{constrained} features, computed by \texttt{constrained\_indexes}); the remaining
indexes $F=\{1,\dots,d\}\setminus K$ are the \emph{free} features.
For $x\in\R^d$ we write $x=(x_K,x_F)$.

\begin{definition}[Constrained region and real profiles]
The constrained region is
\[
  C=\{c\in\R^{K} : c \text{ satisfies } \varphi \text{ and the boundary constraints}\},
\]
which for breastcancer is a box with nonempty interior.
The set of real profiles is the finite set
\[
  P=\{(x_i)_F : 1\le i\le n\},\qquad N=|P|\le n .
\]
The verification domain is $\Dom=C\times P=\{x : x_K\in C,\ x_F\in P\}$.
\end{definition}

Let $\psi$ be the right constraint and $t$ the class it requires.
A learner $L$ maps a training set $D$ to a classifier $f=L(D)$, and $f$ is
\emph{certified} when $f(x)=t$ for every $x\in\Dom$.
In the code, $\Dom$ is enforced by the SMT assertion produced by
\texttt{free\_profile\_domain}, conjoined with $\neg(\varphi\to\psi)$ and the boundary
constraints.

\begin{assumption}[Exact oracle]\label{ass:oracle}
Given $f$, the solver terminates and returns some $x\in\Dom$ with $f(x)\ne t$ whenever
such an $x$ exists, and reports unsatisfiability otherwise.
\end{assumption}

This holds for the forests encoded in linear real arithmetic, a decidable theory for which
Z3 is a complete decision procedure on the generated formulas.

\paragraph{The repair loop.}
Fix an integer $m\ge 1$ (\texttt{profiles\_per\_cycle}, $m=5$ in the experiments).
For $k=0,1,2,\dots$:
\begin{enumerate}
  \item train $f_k=L(D_k)$;
  \item if $f_k$ is certified, stop; otherwise the oracle returns a counterexample
        $x^k=(c^k,p^k)\in\Dom$ with $f_k(x^k)\ne t$;
  \item choose $Q_k\subseteq P\setminus\{p^k\}$ with $|Q_k|=\min(m-1,N-1)$
        (the nearest profiles to $p^k$, see \texttt{nearest\_profiles}) and add the batch
        \[
          A_k=\{(c^k,p^k)\}\cup\{(c^k,q) : q\in Q_k\},
        \]
        every element labelled $t$: $D_{k+1}=D_k\cup(A_k\times\{t\})$.
\end{enumerate}
Each batch adds at most $m$ instances.
We call \emph{repaired} the profiles used by some earlier batch:
\[
  R_k=\{p\in P : (c,p)\in A_j \text{ for some } j<k \text{ and some } c\},\qquad R_0=\emptyset .
\]

\section{The bound}

\begin{assumption}[Slice persistence]\label{ass:persistence}
For every cycle $k$ and every repaired profile $p\in R_k$,
\[
  f_k(c,p)=t\quad\text{for every } c\in C .
\]
In words: once a profile has been moved into the constrained region with label $t$,
every later retrained model classifies the whole slice $C\times\{p\}$ as $t$.
\end{assumption}

\begin{theorem}\label{thm:bound}
Under Assumptions~\ref{ass:oracle} and~\ref{ass:persistence}, the repair loop stops with a
certified model after a number of cycles
\[
  K\le N,
\]
and the number of added instances is
\[
  I=\sum_{k<K}|A_k|\le mK\le mN .
\]
\end{theorem}

\begin{proof}
We first show that every counterexample uses a profile that has not been repaired yet.
Let $f_k$ be not certified and let $x^k=(c^k,p^k)$ be the counterexample returned by the
oracle, so $c^k\in C$ and $f_k(c^k,p^k)\ne t$.
If $p^k$ belonged to $R_k$, Assumption~\ref{ass:persistence} would give
$f_k(c^k,p^k)=t$, a contradiction.
Hence $p^k\notin R_k$.

Since $(c^k,p^k)\in A_k$, we have $p^k\in R_{k+1}$, and therefore
$R_{k+1}\supseteq R_k\cup\{p^k\}$ with $p^k\notin R_k$, that is,
\[
  |R_{k+1}|\ge |R_k|+1 .
\]
By induction from $|R_0|=0$, if the loop has not stopped before cycle $k$ then
$|R_k|\ge k$.
Since $R_k\subseteq P$, we get $k\le|R_k|\le N$.

Suppose the loop reaches cycle $k=N$.
Then $|R_N|=N$, so $R_N=P$, and Assumption~\ref{ass:persistence} gives $f_N(c,p)=t$ for
every $c\in C$ and every $p\in P$, i.e.\ for every $x\in\Dom$: $f_N$ is certified.
By Assumption~\ref{ass:oracle} the oracle then reports unsatisfiability and the loop
stops.
Hence the loop stops at some cycle $K\le N$, and since each batch has at most $m$
elements, $I\le mK\le mN$.
\end{proof}

\begin{corollary}[Fresh neighbours]\label{cor:fresh}
If moreover every batch only uses profiles that are not repaired yet, i.e.\
$Q_k\cap R_k=\emptyset$ for every $k$, then
\[
  K\le\left\lceil N/m\right\rceil
  \qquad\text{and}\qquad
  I\le m\left\lceil N/m\right\rceil<N+m .
\]
\end{corollary}

\begin{proof}
The profiles of $A_k$ are $p^k$ and the $m-1$ elements of $Q_k$, which are pairwise
distinct because $Q_k\subseteq P\setminus\{p^k\}$.
None of them is in $R_k$: $p^k\notin R_k$ by the proof of Theorem~\ref{thm:bound}, and
$Q_k\cap R_k=\emptyset$ by hypothesis.
Hence $|R_{k+1}|=|R_k|+m$, and $|R_k|=mk$ for every cycle $k$ the loop reaches.

The loop produces a counterexample at cycles $0,\dots,K-1$.
At cycle $K-1$ the counterexample profile is not repaired, so $R_{K-1}\ne P$ and
$m(K-1)=|R_{K-1}|<N$.
Therefore $K<N/m+1$, i.e.\ $K\le\lceil N/m\rceil$, and $I=mK\le m\lceil N/m\rceil<N+m$.
\end{proof}

\begin{remark}[Relapses]\label{rem:relapse}
Without any assumption on the learner, counting gives $K\le\rho N$, where $\rho$ is the
largest number of times a single profile is returned as the profile of a counterexample.
Assumption~\ref{ass:persistence} is exactly what makes $\rho=1$.
A counterexample whose profile was already repaired is a \emph{relapse}; every other cycle
repairs a new profile, so at most $N$ cycles are not relapses and $K\le N+(\text{number of
relapses})$.
Section~\ref{sec:experiments} measures how often this happens.
\end{remark}

\begin{remark}[Batches of perturbed copies]
If each cycle also adds $b$ perturbed copies of the counterexample (\texttt{batch\_size}),
the batch has $b+m-1$ elements and the same proof gives $I\le(b+m-1)N$.
Perturbed copies must keep their free part in $P$ for Assumption~\ref{ass:persistence} to
remain meaningful; the current \texttt{perturb} also perturbs the free features, so the
experiments use $b=1$.
\end{remark}

\section{The assumption cannot be dropped}

\begin{proposition}\label{prop:nobound}
There is a learner that fits its training set exactly and an exact oracle for which the
repair loop never stops, even with $N=1$.
Hence, without Assumption~\ref{ass:persistence}, no bound on $K$ or $I$ depending only on
$N$ and $m$ exists.
\end{proposition}

\begin{proof}
Let $s\ne t$ be another class, let $\|\cdot\|$ be the Euclidean norm, and define
\[
  L(D)=f_D,\qquad
  f_D(x)=\begin{cases}
    t & \text{if } \|x-z\|<2^{-|D|} \text{ for some } (z,t)\in D,\\
    s & \text{otherwise.}
  \end{cases}
\]
On any training set whose $t$-labelled and $s$-labelled points are at distance at least
$2^{-|D|}$ from each other, $f_D$ predicts every training label correctly, so the
learner fits its data.

Take $P=\{p\}$ and let $C$ have nonempty interior, so its Lebesgue measure $\lambda(C)$ in
$\R^{K}$ is positive.
At cycle $k$, the set of points of the slice $C\times\{p\}$ on which $f_k$ predicts $t$ is
contained in the union of $|D_k|$ balls of radius $2^{-|D_k|}$, whose total measure in
$\R^{K}$ is at most
\[
  |D_k|\,\omega_{|K|}\,2^{-|K|\,|D_k|},
\]
where $\omega_{|K|}$ is the volume of the unit ball.
This quantity tends to $0$ as $|D_k|$ grows, so for $|D_0|$ large enough it is smaller
than $\lambda(C)$ at every cycle, since $|D_k|$ only increases.
Hence there is always some $c\in C$ with $f_k(c,p)=s\ne t$, the exact oracle always
returns a counterexample, and the loop never stops.
\end{proof}

The learner of Proposition~\ref{prop:nobound} is artificial, but it isolates the mechanism:
a model that repairs only a small neighbourhood of each added instance leaves the rest of
the continuous region $C$ unrepaired.
A decision forest with a minimum leaf size larger than one is the opposite extreme, since a
leaf covers a whole box; this is why one added instance alone rarely changes the forest,
and why the repair adds the nearest profiles too.

\section{Why the original repair has no such bound}

In the original repair the free features range over the box
$B_F=\prod_{j\in F}[\min_i x_{ij},\max_i x_{ij}]$, so the verification domain is
$C\times B_F$ and the role of $P$ is played by the continuum $B_F$.
Even under the analogue of Assumption~\ref{ass:persistence} for these profiles, the proof of
Theorem~\ref{thm:bound} only shows that every counterexample uses a profile not repaired yet;
now there are infinitely many profiles, so the growth of $R_k$ gives
no bound at all.
This matches what we observed: the solver returns profiles at the corners of $B_F$ (most
free features at their minimum), each repair covers a negligible part of $B_F$, and
certifying needed up to 824 instances on breastcancer.

\section{Comparison with the experiments}\label{sec:experiments}

For breastcancer ($n=455$ training instances, $m=5$) the profiles are all distinct, so
$N=455$ for every constraint size and Theorem~\ref{thm:bound} gives $K\le 455$ and
$I\le 2275$.
Table~\ref{tab:observed} reports what the loop actually needed, together with an empirical
check of Assumption~\ref{ass:persistence}: a \emph{relapse} is a counterexample whose profile had already been repaired (as a
counterexample or as a neighbour), which is exactly what the assumption forbids.

\begin{table}[ht]
\centering
\begin{tabular}{llrrrrr}
\toprule
Constrained & Seed & Cycles $K$ & Instances $I$ & Distinct profiles & Relapses & Bound $mN$\\
\midrule
$V_1$ & 1 & 34 & 170 & 30 & 5 & 2275\\
$V_1$ & 2 & 30 & 150 & 26 & 10 & 2275\\
$V_1$ & 3 & 34 & 170 & 29 & 6 & 2275\\
$V_1$ & 4 & 40 & 200 & 35 & 11 & 2275\\
$V_1$ & 5 & 19 & 95 & 18 & 1 & 2275\\
\midrule
$V_1,V_3$ & 1 & 9 & 45 & 9 & 0 & 2275\\
$V_1,V_3$ & 2 & 6 & 30 & 5 & 1 & 2275\\
$V_1,V_3$ & 3 & 6 & 30 & 5 & 1 & 2275\\
$V_1,V_3$ & 4 & 8 & 40 & 8 & 0 & 2275\\
$V_1,V_3$ & 5 & 7 & 35 & 7 & 0 & 2275\\
\midrule
$V_1,V_3,V_4$ & 1 & 5 & 25 & 5 & 0 & 2275\\
$V_1,V_3,V_4$ & 2 & 4 & 20 & 4 & 1 & 2275\\
$V_1,V_3,V_4$ & 3 & 7 & 35 & 7 & 0 & 2275\\
$V_1,V_3,V_4$ & 4 & 5 & 25 & 5 & 0 & 2275\\
$V_1,V_3,V_4$ & 5 & 4 & 20 & 4 & 0 & 2275\\
\bottomrule
\end{tabular}
\caption{Real-profile repair on breastcancer, 8 trees, 5 seeds per constraint size.}
\label{tab:observed}
\end{table}

Assumption~\ref{ass:persistence} held exactly in 7 of the 15 runs, all with two or three constrained features. Over all runs, 36 of the 218 counterexamples (16.5\%) were relapses, almost all of them with only $V_1$ constrained, where the slice $C\times\{p\}$ is widest and a retrained forest can still split inside it. So the assumption is an idealisation and Theorem~\ref{thm:bound} is a bound for an idealised learner; by Remark~\ref{rem:relapse} each relapse adds at most one cycle to the bound, and in every run the observed number of cycles stayed far below the bound: at most 40 cycles and 200 instances, against $N=455$ and $mN=2275$.

\end{document}
